\section{Transporting the inherited loop response}
\label{sec:inherited-loop-transport}
The loop term entering a stationary step already has zero total
coefficient and zero first displacement moments. These cancellations
make its quadratic momentum picture vanish to fourth order. We
prove that the stationary pullback retains this order and gains
$L^{-2}$ in its strip bound, where $L$ is the blocking side.
The order statement and the quantitative strip statement require
different arguments.

\subsection{The fine-lattice picture}
Group the loop coefficients by the two plaquette directions, the
two colour indices, and the relative displacement $z$ of their
bases. Denote the resulting real coefficient array by $r(z)$.
Its exponential norm is bounded by the loop budget $\Lambda$:
the norm of the relative displacement is no greater than the sum
of the two slot lengths. Absolute summability permits grouping
all atoms with the same displacement. The path's moment identities
give
\[
                         \sum_zr(z)=0,\qquad \sum_z z_a r(z)=0.
\]
Consequently its Fourier series $\widehat r(p)$ vanishes to order
two at zero. The quadratic coefficient of the loop action is the
sum over face indices of
\begin{equation}\label{eq:inherited-fine-symbol}
                  d_1(-p)^{\mathsf T}\widehat r(p)d_1(p),
\end{equation}
with its two field arguments symmetrized and its lattice
orientations averaged. Each curl vanishes to first order.
Thus every component of the resulting common picture $S$ has
order four. Its strip bound is at most $C_\tau\Lambda$ at the
original path width $\tau>0$: each curl entry costs at most
$4(e^\tau+1)$, there are finitely many face pairs, and the signed
orientation average has the fixed reflection cost already estimated
in Proposition~\ref{prop:ex-atom-stream-picture}.

Formula \eqref{eq:inherited-fine-symbol} records the actual
finite-volume quadratic coefficient. At zero total momentum, the
two slot offsets appear only through their relative displacement;
the anchor sum gives the number of fine sites. At nonzero total
momentum it is zero. This verifies the identity before taking a
large-volume limit. Invariance and the finite Ward identity give
\begin{equation}\label{eq:inherited-transverse}
                  S(p)d_0(p)=0,\qquad d_0(-p)^{\mathsf T}S(p)=0.
\end{equation}
All tensors below use this unconjugated bilinear convention. On
real momenta, $d_0(-p)^{\mathsf T}=d_0(p)^*$.

\subsection{A strip bound with four finite differences}
Put $\mu=L^{-1}$. Let $X$ be the derivative of the stationary
chart, represented by a common coarse-to-fine Fourier bridge.
Assume that, on a fixed positive coarse strip of width $\sigma$,
its fine differences of order $j\le2$ obey the bounds
\begin{equation}\label{eq:inherited-bridge-price}
                      |\Delta^j X|\le D\mu^{1+j}.
\end{equation}
These are the bridge bounds supplied by the stationary chart.

Apply the finite-difference division of
Lemma~\ref{lem:ex-soft-factorization} to the order-four picture
$S$. It writes $S$ as a finite sum of products of four symbols
$e^{ip_a}-1$ times Fourier series having a common exponential
budget $C\Lambda$ on a positive width $w$ depending on $\tau$.
Distribute two differences to each of the two field arguments.
Moving a difference from momentum $p$ to $-p$ uses
\[
                      q_a(p)=-e^{ip_a}q_a(-p).
\]
The extra phase is a unit lattice translation of the remaining
coefficient; its cost is absorbed in $C$. Thus every term has
exactly two differences on each bridge leg.

Pull this representation back by $X$. On a smaller fixed coarse
strip, choose it narrow enough that all fine phases lie in the
coefficient strip, in particular $\mu\sigma\le w$. The
exponentially weighted coefficient sum bounds every displaced
contraction uniformly. For each pair of legs,
\eqref{eq:inherited-bridge-price} contributes
\[
                         D^2\mu^3\mu^3=D^2\mu^6.
\]
Summing the fine anchor within one coarse block contributes
$L^4=\mu^{-4}$. Hence the complete common coarse picture $H$
has, on a positive width $\sigma'$ depending only on $\sigma$,
component bound
\begin{equation}\label{eq:inherited-strip-gain}
                              |H|_{\sigma'}\le C'\Lambda\mu^2.
\end{equation}
The finite numbers of colour, direction, and difference addresses
are included in $C'$, which depends on the fixed widths and $D$,
not on $L$, the torus, or $\Lambda$. Normal convergence of the
coefficient sums proves holomorphy. Their covariance under
integer momentum shifts and that of $X$ prove periodicity of the
common picture. Expanding the same absolutely convergent sums
on a finite torus gives its exact recording identity for the
linear pullback of the fine loop action.

The nonlinear stationary pullback has the same second derivative.
Indeed the fine loop action $V$ has $DV(0)=0$, and the stationary
map $\psi$ satisfies $\psi(0)=0$ and $D\psi(0)=X$. The chain rule
therefore reads
\[
 D^2(V\circ\psi)(0)[a,b]
       =D^2V(0)[Xa,Xb]+DV(0)[D^2\psi(0)[a,b]]
       =D^2V(0)[Xa,Xb].
\]
The coordinate and tree-gauge maps identify these with the actual
inherited response; changes of allowed tags agree as germs.
Thus $H$ records that response, with no additional error.

\subsection{The alias formula and the coarse--fine equation}
To prove the order at zero, use the real-momentum alias
representation. For $l\in(\mathbb Z/L\mathbb Z)^4$, let
\[
                    k_l(p)=\frac{p+2\pi l}{L}
\]
with any fixed integer representatives; periodicity removes the
choice of representatives. Let $N_l(p)$ be the alias component
of $X$ in the normalization for which a coarse unit wave pulls
back into fine waves of amplitudes $N_l(p)$. With the source
normalization of these amplitudes the common picture satisfies
\begin{equation}\label{eq:inherited-alias}
 H_{ab}(p)=L^4\sum_l
     N_{-l}(-p)_a^{\mathsf T}S(k_l(p))N_l(p)_b.
\end{equation}
For clarity, this equality follows first at each finite coarse
Fourier grid. Expand both pulled-back waves in the finite set
of aliases. Fine translation invariance kills every pair except
opposite aliases; the ratio of fine to coarse site counts is
$L^4$. Approximate an arbitrary real $p$ along grids whose side
lengths tend to infinity. The bridge and both common pictures
are continuous and the alias set is finite, so the identity
passes to every real $p$. The same argument works with
per-volume errors tending to zero. Here both recording identities
are exact.

Write $K_f(k)=d_1(k)^*B(k)d_1(k)$ for the fine Gaussian form and
$K_c(p)$ for the coarse one. The stationary Hessian equations
imply, in the block Fourier coordinates,
\begin{equation}\label{eq:inherited-primal}
          K_f(k_l(p))N_l(p)_a=A_l(p)K_c(p)e_a,
\end{equation}
where $A_l(p)$ is the corresponding component of the adjoint
block-mean map and the inverse block Fourier change of coordinates.
It is analytic near zero for fixed $L$. To obtain
\eqref{eq:inherited-primal}, differentiate the stationary equation,
which gives $K_f X=Q^*K_c$, and diagonalize the fine Gaussian form
in its alias blocks. This identity also follows on arbitrary real
momenta from its exact finite-grid versions by the preceding
density argument. Since $K_c$ contains a curl on each side,
\[
                             K_c(p)=O(|p|^2).
\]

\subsection{Suppressing the nonprincipal aliases}
Fix $l\ne0$ and real $p$ near zero. Then $k=k_l(p)$ stays away
from zero modulo $2\pi$. Let $D(k)=d_0(k)$ and
$\omega(k)=\sum_a|q_a(k)|^2>0$. Project the bridge component
off its longitudinal part:
\[
 N_l^{\perp}(p)
    =\left(I-D(k)\omega(k)^{-1}D(k)^*\right)N_l(p).
\]
The matrix $D(k)^*D(k)$ is $\omega(k)I$ on colours, so
$D(k)^*N_l^{\perp}(p)=0$. The fine kinetic form kills pure
gauge vectors. Combining this fact with
\eqref{eq:inherited-primal} gives
\begin{equation}\label{eq:inherited-elliptic-alias}
     \bigl(K_f(k)+D(k)D(k)^*\bigr)N_l^{\perp}(p)
                           =A_l(p)K_c(p).
\end{equation}

The matrix on the left is invertible at $p=0$. Here the fine
floor is required for the same bed on every finite period, with
one positive lower constant. Passing its Fourier inequality to
real frequencies gives
\[
       \langle v,K_f(k)v\rangle\ge c\|d_1(k)v\|^2,
                                  \qquad c>0.
\]
The coordinate exterior identity is
\[
        \|d_1(k)v\|^2+\|D(k)^*v\|^2=\omega(k)\|v\|^2.
\]
It follows either by expanding both squares or from the
anticommutation of exterior multiplication and contraction.
A vector in the kernel of $K_f(k)+D(k)D(k)^*$ must make both
nonnegative terms zero, and hence must be zero. Finite
dimensionality gives invertibility. The matrix depends
continuously on $p$, so its inverse is bounded near zero for
this fixed $l,L$. Equation~\eqref{eq:inherited-elliptic-alias}
now proves
\[
                              N_l^{\perp}(p)=O(|p|^2).
\]
This is a qualitative order statement. No uniform lower bound
in $L$ is asserted for this inverse.

The Ward identities \eqref{eq:inherited-transverse} permit
replacing both $N$ factors in each alias term of
\eqref{eq:inherited-alias} by their transverse projections.
For $l\ne0$, each contributes order two; $S(k_l(p))$ is bounded
near $p=0$. Every nonprincipal term therefore has order four.
For the principal alias, $k_0(p)=\mu p$, and $S(\mu p)$ already
has order four. Its two bridge factors are analytic and bounded
near zero. This term has order four as well. There are finitely
many aliases, and multiplication by the fixed factor $L^4$
does not change Taylor order. Hence $H$ has order four.

Together with \eqref{eq:inherited-strip-gain} and the exact
recording identity, this supplies the inherited loop picture
used in Section~\ref{sec:stationary-output-construction}.
Its strip bound is uniform in the blocking side under the stated
bridge prices; its order-four property uses the all-period floor
and the exact coarse--fine stationary equation. In particular it
does not assume that the longitudinal alias amplitudes themselves
are small.
